\chapter{Funciones Avanzadas}
\label{ch:advanced-features}

El firmware \textbf{ESP32\_BT\_Controller} está diseñado como una plataforma modular de RC + telemetría.
Más allá del control básico por joysticks, incluye un protocolo de comunicación liviano, modos de proyecto
configurables, un sistema de telemetría estructurado y expansión I2C opcional.

Este capítulo destaca las funciones avanzadas más útiles y explica cómo funcionan en la práctica.

\section{Resumen del Protocolo: Comandos, Eventos y Telemetría}
\label{sec:protocol-overview}

La comunicación se basa en paquetes (ver \texttt{packets.h}). El firmware usa:

\begin{itemize}
    \item \textbf{Paquetes RC State:} estado de control continuo (sticks, perillas, interruptores).
    \item \textbf{Paquetes de Evento:} comandos cortos de una sola vez (pulsaciones de botones).
    \item \textbf{Paquetes de Telemetría:} retroalimentación desde el ESP32 hacia la app Android (paneles, indicadores, gráficas y etiquetas de configuración).
\end{itemize}

Todos los paquetes usan:
\begin{itemize}
    \item Encabezados fijos (p.\ ej., \texttt{0xCC} seguido por un byte de tipo de paquete)
    \item Estructuras empaquetadas para un parseo predecible
    \item Un checksum aditivo simple (liviano y rápido para sistemas embebidos)
\end{itemize}

\section{Sistema de Telemetría (Paneles, Indicadores, Gráficas)}
\label{sec:telemetry}

La telemetría es una de las partes más potentes de este firmware. Permite que el ESP32 envíe
valores en tiempo real de vuelta a la app Android para que el operador pueda \emph{ver lo que el robot está haciendo}
mientras lo conduce.

\subsection{Categorías de telemetría}
El firmware genera múltiples flujos de telemetría (implementados en \texttt{telemetry.cpp}):

\begin{itemize}
    \item \textbf{Paneles} (valores numéricos izquierda/derecha): lecturas compactas como distancia, velocidad, temperatura o cualquier número derivado de un sensor.
    \item \textbf{Indicadores}: un conjunto pequeño de ``widgets de estado'' que incluye:
    \begin{itemize}
        \item un valor analógico (0--100)
        \item un nivel de batería (0--100)
        \item una máscara digital (campo de bits) para múltiples indicadores on/off
    \end{itemize}
    \item \textbf{Gráficas (plots)}: series de actualización rápida para gráficos en el tiempo (las etiquetas por defecto incluyen \texttt{Volts}, \texttt{Amps} y \texttt{RPMs}).
\end{itemize}

\subsection{Frecuencias de actualización de telemetría (significado práctico)}
La telemetría se envía de forma periódica (constantes de tiempo en \texttt{telemetry.cpp}):

\begin{itemize}
    \item \textbf{Gráficas (plots):} cada 50 ms (\(\approx 20\,Hz\)) --- gráficos suaves para señales que cambian rápido.
    \item \textbf{Indicadores:} cada 500 ms (\(\approx 2\,Hz\)) --- valores de UI estables para batería/estado.
\end{itemize}

\subsection{Handshake automático de configuración (etiquetas)}
Cuando el teléfono se conecta, el ESP32 envía un \textbf{paquete de telemetría de configuración una vez por conexión}
(\texttt{sendConfigTelemetry()}). Este paquete proporciona a la app etiquetas legibles para:

\begin{itemize}
    \item nombres de series de gráficas (p.\ ej., \texttt{Volts}, \texttt{Amps}, \texttt{RPMs})
    \item nombres de paneles (p.\ ej., \texttt{Left}, \texttt{Right})
    \item nombres de indicadores (p.\ ej., \texttt{Throttle}, \texttt{Battery})
\end{itemize}

Esta es una característica importante de usabilidad: la app Android puede poblar las etiquetas automáticamente sin
codificarlas de forma fija para cada construcción de robot.

\subsection{Ventana de reenvío de panel (característica de robustez)}
Los valores de panel usan un pequeño mecanismo de reenvío (``ventana de reintento'') de modo que cuando una lectura de panel cambia,
se retransmite repetidamente por un periodo corto (ver \texttt{resendWindow} y \texttt{resendInterval} en
\texttt{telemetry.cpp}). Esto mejora la probabilidad de que el teléfono reciba la lectura actualizada incluso si se pierde un paquete.

\section{Fuentes de Telemetría e Inyección de Debug}
\label{sec:telemetry-sources}

La telemetría está dividida intencionalmente en dos capas:

\begin{itemize}
    \item \textbf{\texttt{telemetry.cpp}} es \emph{solo escritura}: construye paquetes y los envía.
    \item \textbf{\texttt{telemetry\_source.cpp}} decide \emph{de dónde vienen los datos}: lecturas de hardware o valores de depuración.
\end{itemize}

En operación normal, los valores de panel/indicador/gráfica se leen desde el hardware mediante funciones de
\texttt{input\_hw.h} (p.\ ej., entradas numéricas y lecturas ADC). En modos debug, los valores pueden inyectarse por serial
escribiendo valores simples separados por comas. Esto permite verificar el comportamiento de la UI Android y la decodificación
de paquetes \emph{sin} sensores reales conectados.

\paragraph{Por qué esto importa:}
Puedes desarrollar primero la UI + pantallas de telemetría y luego integrar el hardware, reduciendo el tiempo de iteración.

\section{Modos de Proyecto: Solo ESP32 vs Expansión MCP23017}
\label{sec:project-modes}

El firmware puede compilarse en diferentes modos de proyecto (configurados en \texttt{user\_config.h}):

\begin{itemize}
    \item \texttt{MODE\_FULL\_RC}: control RC completo usando GPIO del ESP32 directamente.
    \item \texttt{MODE\_FULL\_RC\_MCP}: control RC completo usando un MCP23017 por I2C para E/S digital expandida.
\end{itemize}

Seleccionar el modo correcto es una ``función avanzada'' porque permite que una sola base de código escale desde robots pequeños
hasta construcciones más grandes con muchos interruptores e indicadores.

\section{Uso del Expansor GPIO MCP23017 (I2C)}
\label{sec:mcp23017}

El MCP23017 agrega 16 pines de E/S digital por I2C y es compatible con el firmware (\texttt{mcp\_io.cpp/.h}).

\subsection{Detalles clave (desde la implementación)}
\begin{itemize}
    \item \textbf{Dirección I2C:} \texttt{0x20}
    \item \textbf{Puerto A (GPIOA):} configurado como \textbf{INPUT} (usado para leer indicadores digitales)
    \item \textbf{Puerto B (GPIOB):} configurado como \textbf{OUTPUT} (usado para salidas de interruptores y eventos)
    \item Las salidas se rastrean por software (\texttt{portBState}) para que los pines individuales puedan alternarse de forma confiable.
\end{itemize}

\subsection{Aplicaciones prácticas}
\begin{itemize}
    \item Una caja de control con muchos interruptores (armar, luces, modos, accesorios)
    \item Un robot con muchas entradas de estado (finales de carrera, líneas de falla, sensores de parachoques)
    \item Arreglos complejos de indicadores LED (sin consumir GPIO del ESP32)
\end{itemize}

\section{Soporte de Sensores I2C (Ejemplos Incluidos)}
\label{sec:i2c-sensors}

El repositorio incluye ejemplos de lectura de sensores I2C en \texttt{i2c\_sensors.cpp}:

\begin{itemize}
    \item Un sensor de temperatura en la dirección \texttt{0x48}
    \item Un ejemplo de IMU en la dirección \texttt{0x68} (lecturas de acelerómetro, p.\ ej., registros estilo MPU6050)
\end{itemize}

Estos ejemplos actúan como plantillas: puedes agregar pasos de inicialización y registros adicionales para soportar
tu módulo de sensor exacto.

\section{Diseño de Entrada Analógica para Compatibilidad con Bluetooth}
\label{sec:adc1}

Una elección de diseño avanzada, sutil pero importante, es el uso de entradas \textbf{solo ADC1} para lecturas analógicas
en la configuración de pines por defecto. En ESP32, ADC2 suele ser problemático cuando Wi-Fi/Bluetooth está activo,
por lo que el mapeo del repositorio mantiene la telemetría analógica en pines con capacidad ADC1 (ver \texttt{pins.h}).

\paragraph{Implicación práctica:}
Si agregas sensores analógicos, prefiere los pines ADC1 configurados (u otros canales ADC1) para evitar lecturas intermitentes
o inválidas mientras Bluetooth está conectado.

\section{Feature Flags (Firmware Modular en Tiempo de Compilación)}
\label{sec:feature-flags}

Las banderas de características en tiempo de compilación se definen en \texttt{feature\_config.h}. Controlan si los subsistemas
se compilan/habilitan (salidas PWM, salidas de interruptores, salidas de eventos, telemetría, funciones de sensores I2C, etc.).

\paragraph{Nota:}
En el snapshot actual del repositorio, \texttt{USE\_TELEMETRY} y \texttt{USE\_I2C\_SENSOR} se definen como
\texttt{0} en ambos modos de proyecto, aunque el código de implementación de telemetría existe en el repositorio. Esto sugiere que la telemetría
puede habilitarse/deshabilitarse por configuración según las necesidades del proyecto (p.\ ej., para reducir ancho de banda o uso de CPU).

\section{Uniendo Todo: Ejemplo de ``Construcción Avanzada''}
\label{sec:advanced-build-example}

Un proyecto avanzado realista combina estas características:

\begin{itemize}
    \item \textbf{Control RC:} dos sticks + perillas para conducción y control de cámara/brazo
    \item \textbf{MCP23017:} muchos interruptores e indicadores usando expansión por I2C
    \item \textbf{Telemetría:} gráficas de batería + corriente, además de paneles para temperatura y distancia
\end{itemize}

Esto crea una experiencia robusta para el operador: el control se mantiene responsivo, mientras que la telemetría ofrece
conciencia situacional y alerta temprana de caída de batería, atascos de motor o fallas de sensores.